\section{Introduction}

The rapid expansion of the mobile app market has resulted in an overwhelming number of new applications being launched each year, yet only a small fraction achieve long-term success. 
Research indicates that approximately 80 -- 90 \% of apps fail within their first year, often due to low user engagement and poor retention \cite{businessOfAppsRetention, ethoraRetention}. 

Even when initial downloads are strong, sustaining user interest remains a major challenge; studies show that up to 77 \% of users stop using an app within the first three days, and the majority abandon it within a month \cite{localyticsRetention}. This highlights a critical issue in digital product design: attracting users is far easier than keeping them engaged. 

Gamification, the use of game design elements such as rewards, progress tracking, and challenges in non-game contexts, has emerged as an effective strategy to address this problem. By leveraging intrinsic and extrinsic motivation, gamification can significantly enhance user engagement, improve retention rates, and ultimately increase the likelihood of an application's success \cite{barariGamification2024}.

To explore the practical application of these concepts, the thesis introduces ``Plant Up!'' \cite{plantUpProject} a gamified smart gardening system developed for this project. By integrating IoT sensors with a mobile application, Plant Up! is designed to transform the routine tracking of plant health into an engaging, interactive experience. It serves as a real-world case study to investigate how digital mechanics can foster sustainable user motivation and bridge the gap between initial interest and long-term engagement.

In light of these objectives, this thesis poses the following primary research question:

\vspace{0.5em}
\noindent
\textit{How can gamification be effectively applied to a gardening system to keep users more engaged in a normal activity?}
\nocite{diplomarbeitSelfDennis}
